\documentclass{svproc}

% Core packages for compact, high-quality Springer svproc typesetting
\usepackage{graphicx}
\usepackage{multirow}
\usepackage{amsmath,amssymb,amsfonts}
\usepackage{mathrsfs}
\usepackage[title]{appendix}
\usepackage{xcolor}
\usepackage{textcomp}
\usepackage{manyfoot}
\usepackage{booktabs}
\usepackage{algorithm}
\usepackage{algorithmicx}
\usepackage{algpseudocode}
\usepackage{subcaption}
\usepackage{placeins}
\usepackage{microtype}
\usepackage{enumitem}
\usepackage{url}
\def\UrlFont{\rmfamily}

% --- Page Budget Control: Strict <= 12 Pages Enforcement ---
\linespread{0.98} 

\makeatletter
% Tighten section heading whitespace
\renewcommand\section{\@startsection{section}{1}{\z@}%
  {-5pt \@plus -1pt \@minus -0.5pt}%
  {2.0pt \@plus 0.5pt}%
  {\normalfont\large\bfseries\boldmath}}
\renewcommand\subsection{\@startsection{subsection}{2}{\z@}%
  {-4.0pt \@plus -0.5pt \@minus -0.5pt}%
  {1.2pt \@plus 0.5pt}%
  {\normalfont\normalsize\bfseries\boldmath}}
\makeatother

\setlist[enumerate]{noitemsep,topsep=0.2pt,parsep=0pt,partopsep=0pt,leftmargin=*}
\setlist[itemize]{noitemsep,topsep=0.2pt,parsep=0pt,partopsep=0pt,leftmargin=*}
\setlength{\textfloatsep}{1.5pt plus 0.5pt minus 0.5pt}
\setlength{\intextsep}{1.0pt plus 0.5pt minus 0.5pt}
\setlength{\floatsep}{1.0pt plus 0.5pt minus 0.5pt}
\setlength{\abovecaptionskip}{0.8pt}
\setlength{\belowcaptionskip}{0.4pt}
\setlength{\abovedisplayskip}{1.0pt plus 0.5pt minus 0.5pt}
\setlength{\belowdisplayskip}{1.0pt plus 0.5pt minus 0.5pt}
\setlength{\abovedisplayshortskip}{0.5pt plus 0.5pt minus 0.5pt}
\setlength{\belowdisplayshortskip}{0.5pt plus 0.5pt minus 0.5pt}
\renewcommand{\topfraction}{0.98}
\renewcommand{\bottomfraction}{0.98}
\renewcommand{\textfraction}{0.02}
\renewcommand{\floatpagefraction}{0.85}

\begin{document}
\mainmatter

\title{OHP-MOCD: Overlapping Community Detection via an Evolutionary Multi-Objective Extension of HP-MOCD}
\titlerunning{OHP-MOCD for Overlapping Community Detection}

\author{Muhammad Anas \and Aneeza Khan \and Manahil Jamil \and Syeda Saleha Raza}
\authorrunning{M. Anas, A. Khan, M. Jamil, and S. S. Raza}

\institute{Department of Computer Science, Habib University, Karachi, Pakistan\\
\url{https://github.com/aanasakhtar/ohp-mocd}}

\maketitle

\begin{abstract}
Overlapping community detection is essential in complex networks where vertices naturally participate in multiple functional modules, yet designing scalable evolutionary multi-objective algorithms (MOEAs) for overlapping structures remains a major challenge. While the recently proposed HP-MOCD framework achieves breakthrough computational efficiency for crisp community detection, its discrete locus-based encoding strictly constrains each vertex to a single cluster. To overcome this limitation, we propose \textbf{OHP-MOCD}, which extends HP-MOCD to overlapping covers via a direct node-to-memberships chromosome. The framework introduces a tri-objective formulation optimized using the Non-dominated Sorting Genetic Algorithm II (NSGA-II), balancing decomposed modularity with an overlap parsimony regularizer. Additionally, OHP-MOCD incorporates an analytically derived, parameter-free statistical boundary bound based on the Central Limit Theorem to govern multi-membership admission. We systematically evaluate OHP-MOCD against six competitive baselines across four distinct algorithmic paradigms on 12 real-world ground-truth networks and synthetic LFR benchmark sweeps ($\mu \in [0.1, 0.6]$). Empirical results demonstrate that OHP-MOCD achieves the highest mean Overlapping NMI on 10 of the 12 real-world networks within the evolutionary paradigm and secures top modularity on dense social topologies, while executing up to $27.5\times$ faster than competing evolutionary methods.
\end{abstract}

\keywords{Overlapping Community Detection, Evolutionary Multi-Objective Optimization, NSGA-II, Modularity Decomposition, Complex Networks, Graph Mining}

\section{Introduction}\label{sec:intro}
Community detection is a foundational problem in network science with broad applications across social network analysis, biological pathway discovery, scientific collaboration modeling, and communication systems \cite{fortunato2010,newman2006,girvan2002}. The objective is to identify functional clusters of vertices characterized by dense internal connectivity and sparse inter-cluster connections. Many traditional approaches formulate this task as a disjoint (hard) partitioning problem, where each vertex belongs to exactly one community. Although this assumption simplifies the combinatorial search space, it is overly restrictive for empirical systems.

In real-world networks, community overlap is a ubiquitous structural property rather than an anomaly \cite{palla2005,xie2013}. A researcher frequently contributes to multiple disciplinary domains, an online social media user participates in distinct familial, professional, and recreational circles \cite{leskovec2012}, and a protein often acts within multiple metabolic complexes. Forcing such boundary vertices into a single cluster obscures critical inter-modular bridges and distorts the underlying network topology. Overlapping community detection addresses this fundamental limitation by allowing vertices to hold simultaneous multi-memberships.

Evolutionary multi-objective algorithms (MOEAs) are particularly well suited to community detection because the structural quality of a community partition cannot be adequately captured by a single scalar metric \cite{pizzuti2012,deb2002}. A high-quality overlapping cover must simultaneously encourage strong intra-community cohesion, minimize inter-community mixing, and prevent spurious overlap bloat. Collapsing these conflicting structural criteria into a single scalar objective requires fixed trade-off hyperparameters that introduce human bias. Multi-objective optimization resolves this by maintaining a Pareto front of non-dominated candidate covers representing diverse topological trade-offs.

Recently, Santos et al.~\cite{santos2025} introduced \textbf{HP-MOCD} (\textit{High-Performance Evolutionary Multi-Objective Community Detection}), achieving breakthrough computational scalability for disjoint community detection by decomposing Newman's modularity into two competing objectives optimized via the Non-dominated Sorting Genetic Algorithm II (NSGA-II) \cite{deb2002}. However, HP-MOCD relies strictly on a discrete Locus-Based Adjacency Representation (LAR) that assigns each vertex to exactly one community label. Consequently, it cannot directly model overlapping covers, even when boundary vertices are structurally supported by multiple groups.

In this work, we propose \textbf{OHP-MOCD}, a principled algorithmic generalization that extends HP-MOCD to overlapping topologies. OHP-MOCD directly maps each vertex to a variable-size membership set containing one primary community and zero or more structurally justified secondary memberships. OHP-MOCD retains the decomposed modularity formulation of HP-MOCD, generalized to overlapping covers via continuous soft membership weights, and introduces a third parsimony objective that regularizes multi-membership complexity. Furthermore, we incorporate Dual-Weighted Interaction (DWI) and an analytically proven, parameter-free statistical boundary bound based on the Central Limit Theorem that admits secondary memberships only when supported by significant neighborhood evidence.

The main contributions of this paper are summarized as follows:
\begin{enumerate}
    \item We generalize HP-MOCD from disjoint partitions to overlapping covers using a direct node-to-memberships chromosome.
    \item We formulate OHP-MOCD as a tri-objective NSGA-II optimization problem, combining decomposed modularity objectives with an overlap complexity regularizer.
    \item We develop topology-aware genetic operators (4-parent ensemble crossover and local-neighborhood majority mutation) tailored for multi-membership chromosomes.
    \item We introduce Dual-Weighted Interaction (DWI) and an analytically derived, parameter-free statistical null boundary bound based on the Central Limit Theorem.
    \item We conduct extensive empirical evaluations across 12 real-world ground-truth networks and synthetic LFR sweeps, benchmarking OHP-MOCD against six baselines across four distinct paradigms.
\end{enumerate}

The remainder of this paper is organized as follows. Section~\ref{sec:related} reviews related work. Section~\ref{sec:problem} formalizes the problem. Section~\ref{sec:method} details the OHP-MOCD methodology and complexity. Section~\ref{sec:experiments} presents the experimental setup, benchmark results, and discussions. Finally, Section~\ref{sec:conclusion} concludes the paper.

\section{Related Work}\label{sec:related}
Community detection was traditionally framed as a disjoint partitioning problem, where Newman and Girvan's modularity $Q$~\cite{newman2004,newman2006} served as the benchmark quality criterion. Palla et al.~\cite{palla2005} demonstrated that overlap in empirical networks is a ubiquitous structural regularity and introduced the Clique Percolation Method. Xie et al.~\cite{xie2013} surveyed overlapping community detection, noting that overlapping vertices rarely exceed 30\% of a network and typically belong to two or three communities---providing an essential parsimony benchmark.

Multi-objective evolutionary search is an effective paradigm for community discovery. Deb et al.'s Non-dominated Sorting Genetic Algorithm II (NSGA-II)~\cite{deb2002} provides the core optimization engine widely adopted across evolutionary graph partitioning methods \cite{pizzuti2012}. Prominent evolutionary algorithms developed for overlapping community detection include MCMOEA by Wen et al.~\cite{wen2016}, which evolves binary indicator strings over maximal clique adjacency graphs; EF-MOCD (Evolutionary Fuzzy Multi-Objective Community Detection) by Tian et al.~\cite{tian2020}, which evolves community centers and derives fuzzy memberships via shortest paths; and MO-EE by Bello-Orgaz et al.~\cite{bello2018}, which evolves line-graph edge partitions.

Non-evolutionary paradigms include the Speaker-Listener Label Propagation Algorithm (SLPA)~\cite{xie2011}, which extends label propagation with dynamic memory retention, and Link Partitioning Around Medoids (LPAM)~\cite{ponomarenko2021}, which performs k-median clustering on commute distances between edges. In neural approaches, NOCD by Shchur and G{\"u}nnemann~\cite{shchur2019} utilizes a 2-layer Graph Convolutional Network under a Bernoulli--Poisson link model.

HP-MOCD by Santos et al.~\cite{santos2025} established a high-performance baseline for disjoint community detection using topology-aware operators and Shi-style modularity decomposition. OHP-MOCD extends this lineage to overlapping covers while preserving its computational efficiency.

\section{Problem Formulation}\label{sec:problem}
Let $G = (V, E)$ be an unweighted, undirected complex network with $N = |V|$ vertices and $M = |E|$ edges. Let $d_v$ denote the degree of vertex $v$. An \textbf{overlapping community cover} is a collection of non-empty subsets $\mathcal{C} = \{C_1, C_2, \dots, C_K\}$ such that $\bigcup_{k=1}^K C_k = V$. Each vertex $v$ is associated with a community membership set $M(v) \subseteq \{1, \dots, K\}$, with overlapping degree $O_v = |M(v)| \ge 1$. When $|M(v)| = 1$ for all $v \in V$, $\mathcal{C}$ reduces to a crisp disjoint partition. We denote $\bar{k} = \frac{1}{N}\sum_{v \in V} |M(v)|$ as the mean community membership count per vertex.

To evaluate partition quality without collapsing multiple conflicting criteria into a single scalar score, OHP-MOCD formulates community discovery as a vector-valued Pareto optimization problem: $\min_{\mathcal{C}} \mathbf{F}(\mathcal{C}) = \big( f_1(\mathcal{C}), f_2(\mathcal{C}), f_3(\mathcal{C}) \Big)$, where $f_1$ measures intra-community soft edge loss, $f_2$ represents the modularity null-model volume penalty, and $f_3$ penalizes excess overlapping memberships.

\section{Proposed Method: OHP-MOCD}\label{sec:method}
OHP-MOCD extends HP-MOCD \cite{santos2025} by replacing the single-label locus representation with a direct node-to-memberships chromosome, introducing an overlap regularization objective, and incorporating topology-aware operators with an analytical boundary rule.

\subsection{Solution Representation}\label{subsec:rep}
Each individual chromosome directly encodes a node-to-memberships mapping: $X = \big(M(v_1), M(v_2), \dots, M(v_N)\big)$, where $M(v_i) = \{c_i^{(1)}, c_i^{(2)}, \dots, c_i^{(k_i)}\}$, with primary community label $c_i^{(1)}$ and optional secondary labels $c_i^{(j)}$ ($j \ge 2$). There is no intermediate decoding step between chromosome and cover; reading a vertex affiliation requires $\mathcal{O}(1)$ operations, and local updates require $\mathcal{O}(\bar{k})$ time.

\subsection{Fitness Function and Objectives}\label{subsec:fitness}
OHP-MOCD decomposes modularity into two competing terms adapted to overlapping covers via uniform soft membership weights $r_{vc} = 1/|M(v)|$. The three fitness objectives are defined as:
\begin{equation}\label{eq:objectives}
f_1 = 1 - \tfrac{1}{M}\!{\textstyle\sum\limits_{(u,v)\in E}}\!{\textstyle\sum\limits_{c}} r_{uc} r_{vc}, \hspace{5pt} f_2 = {\textstyle\sum\limits_{c}}\!\left(\tfrac{\deg(c)}{2M}\right)^{\!2}\!, \hspace{5pt} f_3 = \tfrac{1}{N}\!{\textstyle\sum\limits_{v \in V}} \max(0, |M(v)| - 1),
\end{equation}
where $\deg(c) = \sum_{v:\, c \in M(v)} d_v r_{vc}$ is the soft volume of community $c$. Minimizing $f_1$ maximizes intra-community edge density, minimizing $f_2$ suppresses over-expanded clusters, and minimizing $f_3$ regularizes multi-membership complexity. All three objectives $(f_1, f_2, f_3)$ are optimized jointly throughout the NSGA-II search.

\subsection{Genetic Operators and Environmental Selection}\label{subsec:operators}
OHP-MOCD uses genetic operators that act directly on the node-to-memberships chromosome:
\begin{enumerate}
    \item \textbf{Initialization:} Each vertex $v$ receives a random primary label $c_v^{(1)}$. Under the default boundary-seeded strategy, if the neighborhood $\mathcal{N}(v)$ contains more than one distinct primary community, then with probability $p_{\mathrm{init}} = 0.10$, the strongest runner-up neighboring community is seeded as a secondary membership.
    \item \textbf{Parent Selection:} Four distinct parents are selected via binary tournament selection based on Pareto dominance rank and crowding distance.
    \item \textbf{4-Parent Ensemble Crossover:} With probability $p_c = 0.90$, crossover tallies parent community occurrences: $a_{v,c} = \sum_{i=1}^{4} \mathbb{I}\big(c \in M_i(v)\big)$. The offspring's primary community is the majority vote $c_v^{(1)} \in \arg\max_c a_{v,c}$ (ties broken uniformly at random). Any non-primary community supported by at least two parents ($a_{v,c} \ge 2$) survives as a secondary membership: $M_{\mathrm{child}}(v) = \{c_v^{(1)}\} \cup \{c \neq c_v^{(1)} : a_{v,c} \ge 2\}$.
    \item \textbf{Local-Neighborhood Mutation:} For each vertex selected with probability $p_m = 0.30$, OHP-MOCD counts neighbor community affiliations $m_{v,c} = |\{u \in \mathcal{N}(v) : c \in M(u)\}|$. The primary community is reassigned to the local neighborhood majority $c_v^{(1)} \in \arg\max_c m_{v,c}$, while secondary memberships are updated to retain any community with $m_{v,c} \ge 2$: $M_{\mathrm{mut}}(v) = \{c_v^{(1)}\} \cup \{c \neq c_v^{(1)} : m_{v,c} \ge 2\}$.
\end{enumerate}
Standard NSGA-II fast non-dominated sorting and crowding distance truncation are applied to select the top $P$ survivors from the combined $2P$ population. Following evolution, a memetic boundary refinement step prunes weak secondary memberships where internal support is below $d_v / |M(v)|$.

\subsection{Overlap-Support Mechanisms and Complexity Analysis}\label{subsec:overlap_and_complexity}
\noindent\textbf{1) Dual-Weighted Interaction (DWI):} For vertex $v$ and candidate community $c$, DWI combines topological neighborhood proportion and degree-weighted neighbor influence: $s_{v,c} = \alpha \frac{|\{u \in \mathcal{N}(v) : c \in M(u)\}|}{d_v} + (1-\alpha) \frac{\sum_{u \in \mathcal{N}(v), c \in M(u)} d_u}{\sum_{u \in \mathcal{N}(v)} d_u}$, where $\alpha \in [0, 1]$ balances direct edge consensus and hub influence.

\noindent\textbf{2) Parameter-Free Statistical Null Boundary Bound:} Let $K_v$ be the number of distinct communities adjacent to node $v$. Under the null configuration model where the $d_v$ incident edges are randomly distributed across $K_v$ candidate clusters, the empirical membership proportion follows a multinomial distribution with mean $\mu = \frac{1}{K_v}$ and variance $\sigma^2 \le \frac{1}{4 d_v}$. By the Central Limit Theorem and Chebyshev's inequality, the threshold separating statistically significant multi-membership from stochastic degree noise is bounded by $\theta_v = \frac{1}{K_v} + \frac{1}{\sqrt{d_v}}$. A vertex $v$ is admitted into secondary community $C_c$ if and only if $\frac{|\mathcal{N}(v) \cap C_c|}{d_v} \ge \theta_v$ with $|\mathcal{N}(v) \cap C_c| \ge 2$, eliminating heuristic manual cutoff tuning.

\noindent\textbf{3) Complexity Analysis:} Evaluating an individual costs $\mathcal{O}(N\bar{k} + M\bar{k}_{\deg})$, where $\bar{k} = \frac{1}{N}\sum_u |\mathcal{C}(u)|$ is the mean membership count and $\bar{k}_{\deg} = \frac{1}{M}\sum_{(u,v)\in E} |\mathcal{C}(u)|$ is the degree-weighted overlap. Because overlap remains sparse ($\bar{k} \approx 1$), evaluation scales linearly as $\mathcal{O}(N + M)$ in practice. Four-parent ensemble crossover takes $\mathcal{O}(N)$ per child, and local-neighborhood mutation takes $\mathcal{O}(N + p_m M)$ per child. Both reproduction and fitness evaluation are parallelized across threads via Rayon, while NSGA-II non-dominated sorting and crowding distance require $\mathcal{O}(mP^2)$. Across $G$ generations with $m = 3$ objectives, the total asymptotic runtime complexity is:
\begin{equation}\label{eq:complexity}
\mathcal{O}\Big( G \cdot P(N\bar{k} + M\bar{k}_{\deg} + p_m M) + G \cdot m P^2 \Big).
\end{equation}
Because graph size $(N + M)$ dominates population size $P$, runtime scales with linear edge scans, enabling OHP-MOCD to scale significantly better than baselines reliant on clique enumeration or all-pairs shortest paths.

\section{Experimental Evaluation and Results Discussion}\label{sec:experiments}
Following the empirical structure of Santos et al. \cite{santos2025}, we organize our evaluation into three complementary fronts: (1) an \textbf{intra-paradigm head-to-head comparison} against state-of-the-art MOEAs; (2) a \textbf{cross-paradigm benchmark evaluation} against competitive heuristics, link medoids, and graph neural networks; and (3) a rigorous \textbf{topological modularity, runtime scalability, and synthetic LFR robustness profile}.

\subsection{Experimental Setup, Datasets, and Baseline Suite}\label{subsec:setup}
We benchmark OHP-MOCD against six authentic baseline methods across four major algorithmic paradigms: (1) \textbf{Structural MOEAs:} MCMOEA \cite{wen2016}, EF-MOCD \cite{tian2020}, and MO-EE \cite{bello2018}; (2) \textbf{Dynamic Label Propagation:} SLPA \cite{xie2011}; (3) \textbf{Link Medoids:} LPAM \cite{ponomarenko2021}; (4) \textbf{Graph Neural Networks:} NOCD \cite{shchur2019}. Evaluation metrics include Overlapping Normalized Mutual Information ($ONMI$) \cite{lancichinetti2009onmi}, Pairwise $F_1$-Score, and Shen's Extended Modularity ($EQ$) \cite{shen2009,nicosia2009}. All evolutionary algorithms were evaluated under a common nominal search budget of population size $N_{\text{pop}} = 100$ and $T_{\text{max}} = 100$ generations. Stochastic algorithms were evaluated across 15 independent seeds for small-to-medium networks and 10 seeds for large networks (FB 1684, FB 1912, DBLP). All executions were performed on an identical hardware environment (Intel Xeon / AMD Ryzen 8-core, 32GB RAM). Full source code, baseline implementations, and reproduction scripts are publicly available at \url{https://github.com/aanasakhtar/ohp-mocd}. Topological characteristics of the 12 evaluated real-world ground-truth networks are summarized in Table~\ref{tab:datasets}.

\begin{table}[t]
\centering
\caption{Topological characteristics of the 12 ground-truth benchmark networks.}
\label{tab:datasets}
\resizebox{0.78\textwidth}{!}{%
\begin{tabular}{@{}lrrrrr@{}}
\toprule
\textbf{Network} & \textbf{Nodes ($|V|$)} & \textbf{Edges ($|E|$)} & \textbf{Avg. Degree ($\langle d \rangle$)} & \textbf{Density} & \textbf{Communities ($K$)} \\
\midrule
Karate Club \cite{zachary1977} & 34 & 78 & 4.59 & 0.1390 & 2 \\
Dolphins \cite{lusseau2003} & 62 & 159 & 5.13 & 0.0841 & 2 \\
Polbooks \cite{newman2006} & 105 & 441 & 8.40 & 0.0808 & 3 \\
Football \cite{girvan2002} & 115 & 613 & 10.66 & 0.0935 & 12 \\
Eu-core \cite{leskovec2007} & 1,005 & 16,706 & 33.25 & 0.0331 & 42 \\
Facebook 698 \cite{leskovec2012} & 61 & 270 & 8.85 & 0.1475 & 13 \\
Facebook 414 \cite{leskovec2012} & 150 & 1,693 & 22.57 & 0.1515 & 24 \\
Facebook 686 \cite{leskovec2012} & 170 & 1,656 & 19.48 & 0.1153 & 14 \\
Facebook 348 \cite{leskovec2012} & 224 & 3,192 & 28.50 & 0.1278 & 14 \\
Facebook 1684 \cite{leskovec2012} & 786 & 14,024 & 35.68 & 0.0455 & 17 \\
Facebook 1912 \cite{leskovec2012} & 747 & 30,025 & 80.39 & 0.1078 & 39 \\
DBLP (SNAP) \cite{yang2015} & 1,000 & 3,503 & 7.01 & 0.0070 & 574 \\
\bottomrule
\end{tabular}%
}
\end{table}

\subsection{Front 1: Intra-Paradigm Comparison Against MOEA Baselines}\label{subsec:ea_comparison}
Table~\ref{tab:moea_comparison} presents the dedicated intra-paradigm comparison between OHP-MOCD and competing evolutionary algorithms (MCMOEA \cite{wen2016}, EF-MOCD \cite{tian2020}, and MO-EE \cite{bello2018}) under the exact same $10{,}000$ nominal evaluation budget across all 12 networks, structured identically to the comparative presentation in HP-MOCD \cite{santos2025}.

Within the evolutionary paradigm, OHP-MOCD achieves the highest mean Overlapping NMI ($ONMI$) on 10 of the 12 real-world networks and provides substantial runtime reductions relative to competing evolutionary methods. On the \textit{Football} network, OHP-MOCD reaches $ONMI = 0.7176$, corresponding to improvements of $77.9\%$ and $84.2\%$ over MCMOEA ($0.4034$) and EF-MOCD ($0.3897$), respectively. On the Facebook ego-networks (\textit{Facebook 348, 414, 698}), OHP-MOCD improves $ONMI$ over EF-MOCD by $61.8\%$, $48.8\%$, and $10.6\%$, respectively, while outperforming MCMOEA by $96.2\%$, $64.5\%$, and $31.3\%$. Furthermore, on the densest ego-network (\textit{Facebook 1912}), OHP-MOCD maintains a positive modularity ($EQ = 0.0657$) whereas MCMOEA collapses to negative modularity ($EQ = -0.0011$). In computational efficiency, OHP-MOCD is the fastest evolutionary method on 11 of 12 networks (with MCMOEA marginally faster only on the 34-node Karate network), achieving its largest observed speedups over EF-MOCD on \textit{Eu-core} ($27.5\times$) and \textit{Facebook 1912} ($6.7\times$).

\begin{table}[t]
\centering
\renewcommand{\arraystretch}{0.54}
\setlength{\tabcolsep}{2.5pt}
\caption{Comparison of multi-objective evolutionary community detection algorithms on real-world datasets (nominal $10{,}000$ evaluation budget; mean $\pm$ std across 15 seeds for small-medium networks, 10 seeds for large networks; runtime reports mean wall-clock execution time). Boldface indicates the best mean performance among compared evolutionary methods.}
\label{tab:moea_comparison}
\resizebox{0.94\textwidth}{!}{%
\begin{tabular}{@{}llcccc@{}}
\toprule
\textbf{Dataset} & \textbf{Method} & \textbf{ONMI} & \textbf{Pairwise $F_1$} & \textbf{Modularity ($EQ$)} & \textbf{Runtime (s)} \\
\midrule
% --- Karate ---
Karate & MCMOEA (2016) & $0.1615 \pm 0.0262$ & $0.4393 \pm 0.0260$ & $0.1261 \pm 0.0195$ & $\mathbf{0.31\text{s}}$ \\
 & MO-EE (2018) & $0.1794 \pm 0.0233$ & $0.4921 \pm 0.0244$ & $0.0360 \pm 0.0041$ & $0.51\text{s}$ \\
 & EF-MOCD (2020) & $0.2733 \pm 0.0563$ & $0.5871 \pm 0.0554$ & $0.3538 \pm 0.0130$ & $3.73\text{s}$ \\
 & \textbf{OHP-MOCD (Proposed)} & $\mathbf{0.6571 \pm 0.0763}$ & $\mathbf{0.8349 \pm 0.0544}$ & $\mathbf{0.3886 \pm 0.0143}$ & $0.46\text{s}$ \\
\midrule
% --- Dolphins ---
Dolphins & MCMOEA (2016) & $0.0861 \pm 0.0157$ & $0.3045 \pm 0.0148$ & $0.0921 \pm 0.0048$ & $0.46\text{s}$ \\
 & MO-EE (2018) & $0.0914 \pm 0.0108$ & $0.3110 \pm 0.0154$ & $0.0543 \pm 0.0047$ & $0.69\text{s}$ \\
 & EF-MOCD (2020) & $0.1095 \pm 0.0155$ & $0.3486 \pm 0.0289$ & $0.3895 \pm 0.0143$ & $5.79\text{s}$ \\
 & \textbf{OHP-MOCD (Proposed)} & $\mathbf{0.1354 \pm 0.0120}$ & $\mathbf{0.4405 \pm 0.0151}$ & $\mathbf{0.5192 \pm 0.0104}$ & $\mathbf{0.39\text{s}}$ \\
\midrule
% --- Polbooks ---
Polbooks & MCMOEA (2016) & $0.1324 \pm 0.0091$ & $0.3729 \pm 0.0162$ & $0.0474 \pm 0.0020$ & $1.12\text{s}$ \\
 & MO-EE (2018) & $0.1763 \pm 0.0168$ & $0.5877 \pm 0.0228$ & $0.0307 \pm 0.0025$ & $1.41\text{s}$ \\
 & EF-MOCD (2020) & $0.2594 \pm 0.0353$ & $0.6771 \pm 0.0503$ & $0.3458 \pm 0.0236$ & $9.84\text{s}$ \\
 & \textbf{OHP-MOCD (Proposed)} & $\mathbf{0.3595 \pm 0.0418}$ & $\mathbf{0.7629 \pm 0.0546}$ & $\mathbf{0.4942 \pm 0.0246}$ & $\mathbf{0.65\text{s}}$ \\
\midrule
% --- Football ---
Football & MCMOEA (2016) & $0.4034 \pm 0.0202$ & $0.3560 \pm 0.0217$ & $0.0345 \pm 0.0015$ & $1.38\text{s}$ \\
 & MO-EE (2018) & $0.3312 \pm 0.0268$ & $0.2975 \pm 0.0234$ & $0.0262 \pm 0.0033$ & $1.84\text{s}$ \\
 & EF-MOCD (2020) & $0.3897 \pm 0.0458$ & $0.4200 \pm 0.0376$ & $0.3897 \pm 0.0203$ & $10.68\text{s}$ \\
 & \textbf{OHP-MOCD (Proposed)} & $\mathbf{0.7176 \pm 0.0473}$ & $\mathbf{0.7630 \pm 0.0499}$ & $\mathbf{0.5776 \pm 0.0168}$ & $\mathbf{0.78\text{s}}$ \\
\midrule
% --- Eu-core ---
Eu-core & MCMOEA (2016) & $\mathbf{0.1594 \pm 0.0048}$ & $0.0544 \pm 0.0012$ & $\mathbf{0.0219 \pm 0.0016}$ & $61.95\text{s}$ \\
 & MO-EE (2018) & $0.1520 \pm 0.0079$ & $\mathbf{0.0946 \pm 0.0023}$ & $0.0039 \pm 0.0002$ & $79.80\text{s}$ \\
 & EF-MOCD (2020) & $0.0780 \pm 0.0443$ & $0.0901 \pm 0.0016$ & $0.0195 \pm 0.0158$ & $387.46\text{s}$ \\
 & \textbf{OHP-MOCD (Proposed)} & $0.1440 \pm 0.0259$ & $\mathbf{0.0945 \pm 0.0081}$ & $0.0156 \pm 0.0131$ & $\mathbf{14.08\text{s}}$ \\
\midrule
% --- Facebook 698 ---
Facebook 698 & MCMOEA (2016) & $0.3944 \pm 0.0419$ & $0.5383 \pm 0.0509$ & $0.2131 \pm 0.0263$ & $0.48\text{s}$ \\
 & MO-EE (2018) & $0.4210 \pm 0.0171$ & $0.6311 \pm 0.0367$ & $0.0412 \pm 0.0053$ & $1.00\text{s}$ \\
 & EF-MOCD (2020) & $0.4683 \pm 0.0170$ & $0.6228 \pm 0.0317$ & $0.4366 \pm 0.0279$ & $6.03\text{s}$ \\
 & \textbf{OHP-MOCD (Proposed)} & $\mathbf{0.5177 \pm 0.0253}$ & $\mathbf{0.6651 \pm 0.0336}$ & $\mathbf{0.5154 \pm 0.0071}$ & $\mathbf{0.43\text{s}}$ \\
\midrule
% --- Facebook 414 ---
Facebook 414 & MCMOEA (2016) & $0.3129 \pm 0.0062$ & $0.4396 \pm 0.0078$ & $0.0099 \pm 0.0005$ & $38.19\text{s}$ \\
 & MO-EE (2018) & $0.3161 \pm 0.0186$ & $0.6937 \pm 0.0382$ & $0.0142 \pm 0.0056$ & $4.44\text{s}$ \\
 & EF-MOCD (2020) & $0.3459 \pm 0.0223$ & $0.7302 \pm 0.0316$ & $\mathbf{0.4040 \pm 0.0248}$ & $18.28\text{s}$ \\
 & \textbf{OHP-MOCD (Proposed)} & $\mathbf{0.5146 \pm 0.0449}$ & $\mathbf{0.7320 \pm 0.0672}$ & $0.2909 \pm 0.0449$ & $\mathbf{2.23\text{s}}$ \\
\midrule
% --- Facebook 686 ---
Facebook 686 & MCMOEA (2016) & $0.1652 \pm 0.0035$ & $0.5647 \pm 0.0072$ & $0.0124 \pm 0.0021$ & $32.65\text{s}$ \\
 & MO-EE (2018) & $0.1694 \pm 0.0130$ & $0.6967 \pm 0.0186$ & $0.0081 \pm 0.0012$ & $4.16\text{s}$ \\
 & EF-MOCD (2020) & $0.1563 \pm 0.0140$ & $\mathbf{0.7134 \pm 0.0138}$ & $0.1020 \pm 0.0089$ & $18.73\text{s}$ \\
 & \textbf{OHP-MOCD (Proposed)} & $\mathbf{0.2477 \pm 0.0449}$ & $0.7035 \pm 0.0237$ & $\mathbf{0.1023 \pm 0.0219}$ & $\mathbf{1.78\text{s}}$ \\
\midrule
% --- Facebook 348 ---
Facebook 348 & MCMOEA (2016) & $0.1621 \pm 0.0053$ & $0.2486 \pm 0.0108$ & $0.0091 \pm 0.0004$ & $41.00\text{s}$ \\
 & MO-EE (2018) & $0.2137 \pm 0.0121$ & $0.8249 \pm 0.0310$ & $0.0066 \pm 0.0010$ & $8.70\text{s}$ \\
 & EF-MOCD (2020) & $0.1966 \pm 0.0261$ & $0.8390 \pm 0.0358$ & $\mathbf{0.0936 \pm 0.0220}$ & $31.06\text{s}$ \\
 & \textbf{OHP-MOCD (Proposed)} & $\mathbf{0.3182 \pm 0.0269}$ & $\mathbf{0.8712 \pm 0.0490}$ & $0.0646 \pm 0.0166$ & $\mathbf{3.24\text{s}}$ \\
\midrule
% --- Facebook 1684 ---
Facebook 1684 & MCMOEA (2016) & $0.2025 \pm 0.0049$ & $0.1509 \pm 0.0028$ & $0.0043 \pm 0.0004$ & $42.39\text{s}$ \\
 & MO-EE (2018) & $0.2709 \pm 0.0106$ & $0.4542 \pm 0.0204$ & $0.0081 \pm 0.0007$ & $34.18\text{s}$ \\
 & EF-MOCD (2020) & $0.2274 \pm 0.0235$ & $0.4024 \pm 0.0237$ & $\mathbf{0.2595 \pm 0.0266}$ & $126.23\text{s}$ \\
 & \textbf{OHP-MOCD (Proposed)} & $\mathbf{0.4100 \pm 0.0235}$ & $\mathbf{0.4845 \pm 0.0453}$ & $0.0930 \pm 0.0089$ & $\mathbf{13.06\text{s}}$ \\
\midrule
% --- Facebook 1912 ---
Facebook 1912 & MCMOEA (2016) & $0.1888 \pm 0.0019$ & $0.1873 \pm 0.0023$ & $-0.0011 \pm 0.0001$ & $40.50\text{s}$ \\
 & MO-EE (2018) & $0.2245 \pm 0.0145$ & $0.3777 \pm 0.0136$ & $0.0030 \pm 0.0007$ & $67.47\text{s}$ \\
 & EF-MOCD (2020) & $0.2045 \pm 0.0186$ & $\mathbf{0.4606 \pm 0.0193}$ & $\mathbf{0.2466 \pm 0.0298}$ & $182.49\text{s}$ \\
 & \textbf{OHP-MOCD (Proposed)} & $\mathbf{0.3233 \pm 0.0338}$ & $0.4107 \pm 0.0188$ & $0.0657 \pm 0.0132$ & $\mathbf{27.06\text{s}}$ \\
\midrule
% --- DBLP ---
DBLP & MCMOEA (2016) & $0.3239 \pm 0.0022$ & $0.1000 \pm 0.0023$ & $0.1780 \pm 0.0012$ & $8.41\text{s}$ \\
 & MO-EE (2018) & $\mathbf{0.3271 \pm 0.0091}$ & $0.3023 \pm 0.0308$ & $0.0635 \pm 0.0033$ & $7.01\text{s}$ \\
 & EF-MOCD (2020) & $0.1078 \pm 0.0067$ & $\mathbf{0.4419 \pm 0.0214}$ & $0.3996 \pm 0.0135$ & $64.67\text{s}$ \\
 & \textbf{OHP-MOCD (Proposed)} & $0.2472 \pm 0.0213$ & $0.3443 \pm 0.0579$ & $\mathbf{0.5719 \pm 0.0390}$ & $\mathbf{4.76\text{s}}$ \\
\bottomrule
\end{tabular}%
}
\end{table}

\subsection{Front 2: Cross-Paradigm Benchmark Evaluation}\label{subsec:cross_paradigm}
Table~\ref{tab:cross_paradigm_results} evaluates OHP-MOCD against non-evolutionary paradigms across all 12 benchmark networks: SLPA \cite{xie2011}, LPAM \cite{ponomarenko2021}, and NOCD \cite{shchur2019}. Fig.~\ref{fig:real_world_4panel} visualizes ground-truth recovery metrics and topological scalability across all algorithms.

Across algorithmic paradigms, the results demonstrate distinct and complementary strengths rather than universal dominance. SLPA achieves the highest $ONMI$ on several classical and social benchmarks, whereas OHP-MOCD obtains the highest $ONMI$ on \textit{Facebook 414} ($0.5146$) and the highest Pairwise $F_1$-score on \textit{Facebook 348} ($0.8712$), \textit{Facebook 686} ($0.7035$), \textit{Facebook 698} ($0.6651$), and \textit{DBLP} ($0.3443$). Furthermore, OHP-MOCD achieves top Extended Modularity ($EQ$) on classical benchmarks including \textit{Karate Club} ($0.3886$), \textit{Dolphins} ($0.5192$), \textit{Polbooks} ($0.4942$), and \textit{Facebook 698} ($0.5154$). These results indicate that OHP-MOCD is especially effective at discovering dense, coherent modules in complex social ego-networks and collaboration graphs without manual threshold tuning.

\begin{table}[t]
\centering
\renewcommand{\arraystretch}{0.54}
\setlength{\tabcolsep}{2.5pt}
\caption{Cross-paradigm comparison across all 12 ground-truth networks (mean $\pm$ std across 15 seeds for small-medium networks, 10 seeds for large networks; runtime reports mean wall-clock execution time). Boldface indicates the best mean value across all paradigms.}
\label{tab:cross_paradigm_results}
\resizebox{0.94\textwidth}{!}{%
\begin{tabular}{@{}llcccc@{}}
\toprule
\textbf{Dataset} & \textbf{Method (Paradigm)} & \textbf{ONMI} & \textbf{Pairwise $F_1$} & \textbf{Modularity ($EQ$)} & \textbf{Runtime (s)} \\
\midrule
% --- Karate ---
Karate & SLPA (2011) & $\mathbf{0.7291 \pm 0.1511}$ & $\mathbf{0.8777 \pm 0.0828}$ & $0.3410 \pm 0.0961$ & $0.03\text{s}$ \\
 & LPAM (2021) & $0.4551 \pm 0.1146$ & $0.7387 \pm 0.0855$ & $0.2755 \pm 0.0697$ & $\mathbf{0.02\text{s}}$ \\
 & NOCD (2019) & $0.3613 \pm 0.0000$ & $0.5941 \pm 0.0000$ & $-0.0053 \pm 0.0000$ & $0.02\text{s}$ \\
 & \textbf{OHP-MOCD (Proposed)} & $0.6571 \pm 0.0763$ & $0.8349 \pm 0.0544$ & $\mathbf{0.3886 \pm 0.0143}$ & $0.46\text{s}$ \\
\midrule
% --- Dolphins ---
Dolphins & SLPA (2011) & $\mathbf{0.1447 \pm 0.0174}$ & $\mathbf{0.4578 \pm 0.0485}$ & $0.4554 \pm 0.0596$ & $0.07\text{s}$ \\
 & LPAM (2021) & $0.1138 \pm 0.0187$ & $0.4068 \pm 0.0367$ & $0.4143 \pm 0.0501$ & $0.12\text{s}$ \\
 & NOCD (2019) & $0.0745 \pm 0.0004$ & $0.3953 \pm 0.0010$ & $0.2348 \pm 0.0017$ & $\mathbf{0.01\text{s}}$ \\
 & \textbf{OHP-MOCD (Proposed)} & $0.1354 \pm 0.0120$ & $0.4405 \pm 0.0151$ & $\mathbf{0.5192 \pm 0.0104}$ & $0.39\text{s}$ \\
\midrule
% --- Polbooks ---
Polbooks & SLPA (2011) & $\mathbf{0.4039 \pm 0.0557}$ & $\mathbf{0.7892 \pm 0.0439}$ & $0.4829 \pm 0.0254$ & $0.12\text{s}$ \\
 & LPAM (2021) & $0.2553 \pm 0.0314$ & $0.5813 \pm 0.0618$ & $0.3459 \pm 0.0462$ & $0.39\text{s}$ \\
 & NOCD (2019) & $0.2339 \pm 0.0007$ & $0.6139 \pm 0.0005$ & $0.1299 \pm 0.0018$ & $\mathbf{0.01\text{s}}$ \\
 & \textbf{OHP-MOCD (Proposed)} & $0.3595 \pm 0.0418$ & $0.7629 \pm 0.0546$ & $\mathbf{0.4942 \pm 0.0246}$ & $0.65\text{s}$ \\
\midrule
% --- Football ---
Football & SLPA (2011) & $\mathbf{0.7750 \pm 0.0327}$ & $\mathbf{0.8270 \pm 0.0521}$ & $\mathbf{0.5916 \pm 0.0124}$ & $0.15\text{s}$ \\
 & LPAM (2021) & $0.3390 \pm 0.0459$ & $0.4029 \pm 0.0371$ & $0.3703 \pm 0.0265$ & $0.22\text{s}$ \\
 & NOCD (2019) & $0.3484 \pm 0.0006$ & $0.1478 \pm 0.0000$ & $0.1233 \pm 0.0000$ & $\mathbf{0.01\text{s}}$ \\
 & \textbf{OHP-MOCD (Proposed)} & $0.7176 \pm 0.0473$ & $0.7630 \pm 0.0499$ & $0.5776 \pm 0.0168$ & $0.78\text{s}$ \\
\midrule
% --- Eu-core ---
Eu-core & SLPA (2011) & $0.1932 \pm 0.0131$ & $0.0893 \pm 0.0032$ & $-0.0034 \pm 0.0297$ & $6.49\text{s}$ \\
 & LPAM (2021) & $\mathbf{0.2051 \pm 0.0177}$ & $\mathbf{0.1658 \pm 0.0304}$ & $\mathbf{0.1928 \pm 0.0281}$ & $92.78\text{s}$ \\
 & NOCD (2019) & $0.1380 \pm 0.0022$ & $0.1032 \pm 0.0002$ & $0.0334 \pm 0.0009$ & $\mathbf{0.15\text{s}}$ \\
 & \textbf{OHP-MOCD (Proposed)} & $0.1440 \pm 0.0259$ & $0.0945 \pm 0.0081$ & $0.0156 \pm 0.0131$ & $14.08\text{s}$ \\
\midrule
% --- Facebook 698 ---
Facebook 698 & SLPA (2011) & $\mathbf{0.5338 \pm 0.0132}$ & $0.5710 \pm 0.0321$ & $0.4891 \pm 0.0141$ & $0.08\text{s}$ \\
 & LPAM (2021) & $0.3273 \pm 0.0488$ & $0.4893 \pm 0.0845$ & $0.3583 \pm 0.0841$ & $0.07\text{s}$ \\
 & NOCD (2019) & $0.4304 \pm 0.0000$ & $0.5198 \pm 0.0000$ & $0.2652 \pm 0.0000$ & $\mathbf{0.01\text{s}}$ \\
 & \textbf{OHP-MOCD (Proposed)} & $0.5177 \pm 0.0253$ & $\mathbf{0.6651 \pm 0.0336}$ & $\mathbf{0.5154 \pm 0.0071}$ & $0.43\text{s}$ \\
\midrule
% --- Facebook 414 ---
Facebook 414 & SLPA (2011) & $0.4964 \pm 0.0121$ & $0.7857 \pm 0.0000$ & $\mathbf{0.5437 \pm 0.0001}$ & $0.32\text{s}$ \\
 & LPAM (2021) & $0.3743 \pm 0.0348$ & $0.6991 \pm 0.0529$ & $0.4523 \pm 0.0474$ & $1.25\text{s}$ \\
 & NOCD (2019) & $0.3966 \pm 0.0010$ & $\mathbf{0.8164 \pm 0.0014}$ & $0.2379 \pm 0.0018$ & $\mathbf{0.02\text{s}}$ \\
 & \textbf{OHP-MOCD (Proposed)} & $\mathbf{0.5146 \pm 0.0449}$ & $0.7320 \pm 0.0672$ & $0.2909 \pm 0.0449$ & $2.23\text{s}$ \\
\midrule
% --- Facebook 686 ---
Facebook 686 & SLPA (2011) & $\mathbf{0.3204 \pm 0.0411}$ & $0.7007 \pm 0.0428$ & $0.1146 \pm 0.0498$ & $0.33\text{s}$ \\
 & LPAM (2021) & $0.1865 \pm 0.0190$ & $0.3750 \pm 0.0471$ & $\mathbf{0.1887 \pm 0.0321}$ & $2.72\text{s}$ \\
 & NOCD (2019) & $0.1873 \pm 0.0022$ & $0.7000 \pm 0.0006$ & $0.0753 \pm 0.0005$ & $\mathbf{0.02\text{s}}$ \\
 & \textbf{OHP-MOCD (Proposed)} & $0.2477 \pm 0.0449$ & $\mathbf{0.7035 \pm 0.0237}$ & $0.1023 \pm 0.0219$ & $1.78\text{s}$ \\
\midrule
% --- Facebook 348 ---
Facebook 348 & SLPA (2011) & $\mathbf{0.4009 \pm 0.0222}$ & $0.8304 \pm 0.0504$ & $0.1027 \pm 0.0316$ & $0.55\text{s}$ \\
 & LPAM (2021) & $0.2282 \pm 0.0299$ & $0.5728 \pm 0.1478$ & $\mathbf{0.1287 \pm 0.0369}$ & $10.46\text{s}$ \\
 & NOCD (2019) & $0.1891 \pm 0.0011$ & $0.7768 \pm 0.0007$ & $0.0543 \pm 0.0004$ & $\mathbf{0.03\text{s}}$ \\
 & \textbf{OHP-MOCD (Proposed)} & $0.3182 \pm 0.0269$ & $\mathbf{0.8712 \pm 0.0490}$ & $0.0646 \pm 0.0166$ & $3.24\text{s}$ \\
\midrule
% --- Facebook 1684 ---
Facebook 1684 & SLPA (2011) & $\mathbf{0.5019 \pm 0.0312}$ & $0.5386 \pm 0.0349$ & $\mathbf{0.5062 \pm 0.0098}$ & $2.20\text{s}$ \\
 & LPAM (2021) & $0.2864 \pm 0.0182$ & $\mathbf{0.5847 \pm 0.0515}$ & $0.4249 \pm 0.0246$ & $27.01\text{s}$ \\
 & NOCD (2019) & $0.2598 \pm 0.0048$ & $0.5473 \pm 0.0043$ & $0.1136 \pm 0.0027$ & $\mathbf{0.47\text{s}}$ \\
 & \textbf{OHP-MOCD (Proposed)} & $0.4100 \pm 0.0235$ & $0.4845 \pm 0.0453$ & $0.0930 \pm 0.0089$ & $13.06\text{s}$ \\
\midrule
% --- Facebook 1912 ---
Facebook 1912 & SLPA (2011) & $\mathbf{0.4089 \pm 0.0173}$ & $\mathbf{0.6643 \pm 0.0037}$ & $\mathbf{0.5264 \pm 0.0007}$ & $3.93\text{s}$ \\
 & LPAM (2021) & $0.2162 \pm 0.0265$ & $0.5363 \pm 0.0731$ & $0.3547 \pm 0.1376$ & $62.96\text{s}$ \\
 & NOCD (2019) & $0.1891 \pm 0.0017$ & $0.5273 \pm 0.0022$ & $0.0768 \pm 0.0025$ & $\mathbf{0.09\text{s}}$ \\
 & \textbf{OHP-MOCD (Proposed)} & $0.3233 \pm 0.0338$ & $0.4107 \pm 0.0188$ & $0.0657 \pm 0.0132$ & $27.06\text{s}$ \\
\midrule
% --- DBLP ---
DBLP & SLPA (2011) & $\mathbf{0.2884 \pm 0.0089}$ & $0.2776 \pm 0.0829$ & $\mathbf{0.6836 \pm 0.0147}$ & $0.80\text{s}$ \\
 & LPAM (2021) & $0.2140 \pm 0.0124$ & $0.2443 \pm 0.0200$ & $0.6295 \pm 0.0404$ & $10.31\text{s}$ \\
 & NOCD (2019) & $0.2131 \pm 0.0048$ & $0.2779 \pm 0.0021$ & $0.2964 \pm 0.0130$ & $\mathbf{0.08\text{s}}$ \\
 & \textbf{OHP-MOCD (Proposed)} & $0.2472 \pm 0.0213$ & $\mathbf{0.3443 \pm 0.0579}$ & $0.5719 \pm 0.0390$ & $4.76\text{s}$ \\
\bottomrule
\end{tabular}%
}
\end{table}

\subsection{Front 3A: Topological Modularity Quality and Scalability}\label{subsec:modularity_and_runtime}
OHP-MOCD demonstrates competitive topological modularity, obtaining top $EQ$ on \textit{Karate Club} ($EQ = 0.3886$), \textit{Dolphins} ($EQ = 0.5192$), \textit{Polbooks} ($EQ = 0.4942$), and \textit{Facebook 698} ($EQ = 0.5154$). In computational scalability, while lightweight heuristics such as NOCD and SLPA execute in under a second on smaller graphs, OHP-MOCD is substantially faster than competing evolutionary algorithms, scaling to $30{,}000$ edges in $27.06\text{s}$ (compared to $182.49\text{s}$ for EF-MOCD and $67.47\text{s}$ for MO-EE), as visualized in Fig.~\ref{fig:real_world_4panel}.

\begin{figure}[t]
\centering
\begin{subfigure}[b]{0.48\textwidth}
    \centering
    \includegraphics[width=\textwidth,height=4.5cm,keepaspectratio]{fig1_onmi_comparison.png}
    \caption{Overlapping NMI ($ONMI$)}
    \label{fig:real_world_onmi}
\end{subfigure}\hfill
\begin{subfigure}[b]{0.48\textwidth}
    \centering
    \includegraphics[width=\textwidth,height=4.5cm,keepaspectratio]{fig2_f1_comparison.png}
    \caption{Pairwise $F_1$-Score}
    \label{fig:real_world_f1}
\end{subfigure}
\begin{subfigure}[b]{0.48\textwidth}
    \centering
    \includegraphics[width=\textwidth,height=4.5cm,keepaspectratio]{fig3_eq_modularity_comparison.png}
    \caption{Extended Modularity ($EQ$)}
    \label{fig:real_world_eq}
\end{subfigure}\hfill
\begin{subfigure}[b]{0.48\textwidth}
    \centering
    \includegraphics[width=\textwidth,height=4.5cm,keepaspectratio]{fig4_runtime_scalability.png}
    \caption{Runtime Scalability (s, log-scale)}
    \label{fig:real_world_time}
\end{subfigure}
\caption{Comparative evaluation across all 12 real-world benchmark networks: Ground-truth recovery ($ONMI$ and $F_1$, top) alongside topological modularity ($EQ$) and computational runtime scalability (bottom).}
\label{fig:real_world_4panel}
\end{figure}

\subsection{Front 3B: Robustness on Synthetic LFR Benchmark Networks}\label{subsec:lfr_sweeps}
To evaluate resilience against topological noise, we conducted parameter sweeps on synthetic LFR benchmark graphs \cite{lancichinetti2008lfr} ($N = 300$, $\langle d \rangle = 12$, community sizes $\in [20, 50]$) across mixing parameter values $\mu \in [0.1, 0.6]$. Table~\ref{tab:lfr_results} provides the quantitative scorecard, while Fig.~\ref{fig:lfr_curves} depicts the degradation curves for $ONMI$, $F_1$, and $EQ$.

At low-to-moderate noise ($\mu \le 0.3$), OHP-MOCD performs strongly and consistently outperforms all competing evolutionary baselines in $ONMI$ and Pairwise $F_1$, achieving $ONMI = 0.9350$ at $\mu = 0.1$ and $0.5976$ at $\mu = 0.2$ (surpassing EF-MOCD: $0.4120$, MO-EE: $0.2495$, MCMOEA: $0.2206$). Although SLPA maintains higher $ONMI$ across the sweep, OHP-MOCD demonstrates a more stable structural modularity profile, maintaining positive modularity ($EQ \ge 0.1293$) across the full range of $\mu$, whereas SLPA collapses to negative modularity ($EQ < 0$) for $\mu \ge 0.4$.

% --- Full Table 4 (renders directly before Conclusion) ---
\begin{table}[t]
\centering
\fontsize{6.2pt}{7.2pt}\selectfont
\renewcommand{\arraystretch}{0.48}
\setlength{\tabcolsep}{2.0pt}
\caption{Comparative performance on synthetic LFR benchmark networks across mixing parameter values $\mu \in [0.1, 0.6]$ (mean $\pm$ std across 5 independent seeds; runtime reports mean wall-clock execution time). Boldface indicates the best mean value among all compared methods.}
\label{tab:lfr_results}
\resizebox{0.94\textwidth}{!}{%
\begin{tabular}{@{}clcccc@{}}
\toprule
\textbf{Mixing Parameter ($\mu$)} & \textbf{Algorithm} & \textbf{ONMI} & \textbf{Pairwise $F_1$} & \textbf{Modularity ($EQ$)} & \textbf{Runtime (s)} \\
\midrule
$\mu = 0.1$ & EF-MOCD (2020) & $0.4120 \pm 0.0218$ & $0.5487 \pm 0.0243$ & $0.3794 \pm 0.0240$ & $27.03\text{s}$ \\
 & LPAM (2021) & $0.4406 \pm 0.0917$ & $0.6227 \pm 0.1040$ & $0.4109 \pm 0.0634$ & $1.41\text{s}$ \\
 & MCMOEA (2016) & $0.2206 \pm 0.0238$ & $0.3198 \pm 0.0094$ & $0.0146 \pm 0.0008$ & $12.69\text{s}$ \\
 & MO-EE (2018) & $0.2495 \pm 0.0306$ & $0.3967 \pm 0.0244$ & $0.0162 \pm 0.0025$ & $8.23\text{s}$ \\
 & NOCD (2019) & $0.5473 \pm 0.0270$ & $0.4367 \pm 0.0340$ & $0.1764 \pm 0.0356$ & $\mathbf{0.16\text{s}}$ \\
 & \textbf{OHP-MOCD (Proposed)} & $0.9350 \pm 0.0265$ & $0.9710 \pm 0.0285$ & $0.5997 \pm 0.0700$ & $2.15\text{s}$ \\
 & SLPA (2011) & $\mathbf{1.0000 \pm 0.0000}$ & $\mathbf{1.0000 \pm 0.0000}$ & $\mathbf{0.6682 \pm 0.0168}$ & $0.80\text{s}$ \\
\midrule
$\mu = 0.2$ & EF-MOCD (2020) & $0.1702 \pm 0.0171$ & $0.3556 \pm 0.0656$ & $0.1602 \pm 0.0570$ & $35.87\text{s}$ \\
 & LPAM (2021) & $0.2371 \pm 0.0330$ & $0.3804 \pm 0.0321$ & $0.2196 \pm 0.0226$ & $1.18\text{s}$ \\
 & MCMOEA (2016) & $0.1855 \pm 0.0145$ & $0.3043 \pm 0.0208$ & $0.0144 \pm 0.0015$ & $14.53\text{s}$ \\
 & MO-EE (2018) & $0.1533 \pm 0.0189$ & $0.3440 \pm 0.0377$ & $0.0106 \pm 0.0015$ & $7.42\text{s}$ \\
 & NOCD (2019) & $0.4577 \pm 0.0518$ & $0.4090 \pm 0.0460$ & $0.0954 \pm 0.0282$ & $\mathbf{0.02\text{s}}$ \\
 & \textbf{OHP-MOCD (Proposed)} & $0.5976 \pm 0.0837$ & $0.6017 \pm 0.1141$ & $0.3352 \pm 0.0435$ & $1.73\text{s}$ \\
 & SLPA (2011) & $\mathbf{0.8838 \pm 0.0960}$ & $\mathbf{0.8815 \pm 0.1112}$ & $\mathbf{0.4979 \pm 0.0589}$ & $0.69\text{s}$ \\
\midrule
$\mu = 0.3$ & EF-MOCD (2020) & $0.1158 \pm 0.0314$ & $0.3161 \pm 0.0597$ & $0.0786 \pm 0.0160$ & $19.10\text{s}$ \\
 & LPAM (2021) & $0.1082 \pm 0.0136$ & $0.2605 \pm 0.0357$ & $0.1467 \pm 0.0124$ & $0.71\text{s}$ \\
 & MCMOEA (2016) & $0.1433 \pm 0.0216$ & $0.2453 \pm 0.0254$ & $0.0127 \pm 0.0007$ & $12.67\text{s}$ \\
 & MO-EE (2018) & $0.1072 \pm 0.0235$ & $0.3005 \pm 0.0390$ & $0.0090 \pm 0.0007$ & $4.23\text{s}$ \\
 & NOCD (2019) & $0.3471 \pm 0.0654$ & $0.3077 \pm 0.0418$ & $0.0333 \pm 0.0095$ & $\mathbf{0.02\text{s}}$ \\
 & \textbf{OHP-MOCD (Proposed)} & $0.3599 \pm 0.1329$ & $\mathbf{0.4418 \pm 0.1580}$ & $\mathbf{0.1829 \pm 0.0887}$ & $1.56\text{s}$ \\
 & SLPA (2011) & $\mathbf{0.4498 \pm 0.1902}$ & $0.3721 \pm 0.1267$ & $0.0652 \pm 0.1579$ & $0.34\text{s}$ \\
\midrule
$\mu = 0.4$ & EF-MOCD (2020) & $0.0628 \pm 0.0521$ & $\mathbf{0.3474 \pm 0.0727}$ & $0.0385 \pm 0.0177$ & $19.11\text{s}$ \\
 & LPAM (2021) & $0.0779 \pm 0.0359$ & $0.2262 \pm 0.0241$ & $0.1315 \pm 0.0167$ & $1.17\text{s}$ \\
 & MCMOEA (2016) & $0.1084 \pm 0.0236$ & $0.2032 \pm 0.0175$ & $0.0123 \pm 0.0012$ & $12.56\text{s}$ \\
 & MO-EE (2018) & $0.0755 \pm 0.0166$ & $0.3183 \pm 0.0525$ & $0.0084 \pm 0.0009$ & $4.26\text{s}$ \\
 & NOCD (2019) & $0.1843 \pm 0.0546$ & $0.3361 \pm 0.0616$ & $0.0103 \pm 0.0078$ & $\mathbf{0.02\text{s}}$ \\
 & \textbf{OHP-MOCD (Proposed)} & $0.0815 \pm 0.0354$ & $0.2734 \pm 0.0534$ & $\mathbf{0.1564 \pm 0.0338}$ & $1.51\text{s}$ \\
 & SLPA (2011) & $\mathbf{0.4134 \pm 0.1936}$ & $0.3450 \pm 0.0715$ & $-0.0045 \pm 0.0029$ & $0.34\text{s}$ \\
\midrule
$\mu = 0.5$ & EF-MOCD (2020) & $0.0275 \pm 0.0239$ & $0.3024 \pm 0.0288$ & $0.0323 \pm 0.0115$ & $19.30\text{s}$ \\
 & LPAM (2021) & $0.0548 \pm 0.0163$ & $0.1997 \pm 0.0203$ & $0.1239 \pm 0.0128$ & $0.77\text{s}$ \\
 & MCMOEA (2016) & $0.0993 \pm 0.0145$ & $0.1632 \pm 0.0063$ & $0.0117 \pm 0.0003$ & $12.74\text{s}$ \\
 & MO-EE (2018) & $0.0624 \pm 0.0098$ & $0.2699 \pm 0.0316$ & $0.0074 \pm 0.0008$ & $4.31\text{s}$ \\
 & NOCD (2019) & $0.1047 \pm 0.0513$ & $0.2868 \pm 0.0421$ & $0.0066 \pm 0.0024$ & $\mathbf{0.02\text{s}}$ \\
 & \textbf{OHP-MOCD (Proposed)} & $0.0688 \pm 0.0270$ & $0.2398 \pm 0.0466$ & $\mathbf{0.1293 \pm 0.0613}$ & $1.55\text{s}$ \\
 & SLPA (2011) & $\mathbf{0.5000 \pm 0.0000}$ & $\mathbf{0.3056 \pm 0.0324}$ & $-0.0037 \pm 0.0008$ & $0.34\text{s}$ \\
\midrule
$\mu = 0.6$ & EF-MOCD (2020) & $0.0128 \pm 0.0044$ & $0.2856 \pm 0.0483$ & $0.0270 \pm 0.0034$ & $19.57\text{s}$ \\
 & LPAM (2021) & $0.0318 \pm 0.0059$ & $0.1733 \pm 0.0365$ & $0.1090 \pm 0.0149$ & $1.10\text{s}$ \\
 & MCMOEA (2016) & $0.0745 \pm 0.0121$ & $0.1470 \pm 0.0091$ & $0.0116 \pm 0.0006$ & $13.66\text{s}$ \\
 & MO-EE (2018) & $0.0520 \pm 0.0068$ & $0.2624 \pm 0.0402$ & $0.0074 \pm 0.0007$ & $4.38\text{s}$ \\
 & NOCD (2019) & $0.0345 \pm 0.0056$ & $0.2825 \pm 0.0467$ & $0.0039 \pm 0.0032$ & $\mathbf{0.01\text{s}}$ \\
 & \textbf{OHP-MOCD (Proposed)} & $0.0300 \pm 0.0121$ & $0.1607 \pm 0.0193$ & $\mathbf{0.1565 \pm 0.0125}$ & $1.51\text{s}$ \\
 & SLPA (2011) & $\mathbf{0.5000 \pm 0.0000}$ & $\mathbf{0.2879 \pm 0.0487}$ & $-0.0030 \pm 0.0007$ & $0.34\text{s}$ \\
\bottomrule
\end{tabular}%
}
\end{table}

\begin{figure}[t]
\centering
\begin{subfigure}[b]{0.32\textwidth}
    \centering
    \includegraphics[width=\textwidth,height=4.2cm,keepaspectratio]{fig5_lfr_onmi_vs_mu.png}
    \caption{$ONMI$ vs. $\mu$}
    \label{fig:lfr_onmi}
\end{subfigure}\hfill
\begin{subfigure}[b]{0.32\textwidth}
    \centering
    \includegraphics[width=\textwidth,height=4.2cm,keepaspectratio]{fig6_lfr_f1_vs_mu.png}
    \caption{Pairwise $F_1$ vs. $\mu$}
    \label{fig:lfr_f1}
\end{subfigure}\hfill
\begin{subfigure}[b]{0.32\textwidth}
    \centering
    \includegraphics[width=\textwidth,height=4.2cm,keepaspectratio]{fig7_lfr_eq_modularity_vs_mu.png}
    \caption{Modularity $EQ$ vs. $\mu$}
    \label{fig:lfr_eq}
\end{subfigure}
\caption{Synthetic LFR Benchmark: Community recovery metrics ($ONMI$, $F_1$, and $EQ$) as a function of topological mixing parameter $\mu \in [0.1, 0.6]$ across seven evaluated methods (six baselines plus proposed OHP-MOCD).}
\label{fig:lfr_curves}
\end{figure}

\FloatBarrier % Flushes Table 4 and Fig 2 before Conclusion & References

\section{Conclusion}\label{sec:conclusion}
In this paper, we presented \textbf{OHP-MOCD}, a principled algorithmic extension of HP-MOCD \cite{santos2025} for overlapping community detection in complex networks. By combining a direct node-to-memberships chromosome, a tri-objective formulation with overlap parsimony regularization, Dual-Weighted Interaction, and an analytically derived statistical boundary bound based on the Central Limit Theorem, OHP-MOCD eliminates heuristic threshold tuning while preserving high-performance evolutionary search. Comprehensive experiments on 12 real-world ground-truth networks and synthetic LFR sweeps established that OHP-MOCD achieves substantial improvements in ground-truth recovery and modularity quality over existing evolutionary multi-objective methods while delivering up to $27.5\times$ computational speedups.

{\scriptsize
\begin{thebibliography}{99}
\setlength{\itemsep}{-0.5pt}
\setlength{\parsep}{0pt}
\setlength{\parskip}{0pt}

% --- 1. Foundational Predecessors & Evolutionary Optimization ---
\bibitem{fortunato2010}
Fortunato, S.: Community detection in graphs. Physics Reports \textbf{486}(3-5), 75--174 (2010)

\bibitem{newman2006}
Newman, M.E.: Modularity and community structure in networks. Proceedings of the National Academy of Sciences \textbf{103}(23), 8577--8582 (2006)

\bibitem{girvan2002}
Girvan, M., Newman, M.E.: Community structure in social and biological networks. Proceedings of the National Academy of Sciences \textbf{99}(12), 7821--7826 (2002)

\bibitem{palla2005}
Palla, G., Der{\'e}nyi, I., Farkas, I., Vicsek, T.: Uncovering the overlapping community structure of complex networks in nature and society. Nature \textbf{435}(7043), 814--818 (2005)

\bibitem{xie2013}
Xie, J., Kelley, S., Szymanski, B.K.: Overlapping community detection in networks: The state-of-the-art and comparative study. ACM Computing Surveys \textbf{45}(4), 1--35 (2013)

\bibitem{leskovec2012}
McAuley, J., Leskovec, J.: Learning to discover social circles in ego networks. In: Advances in Neural Information Processing Systems (NeurIPS), vol. 25 (2012)

\bibitem{pizzuti2012}
Pizzuti, C.: A multiobjective genetic algorithm to find communities in complex networks. IEEE Transactions on Evolutionary Computation \textbf{16}(3), 418--430 (2012)

\bibitem{deb2002}
Deb, K., Pratap, A., Agarwal, S., Meyarivan, T.: A fast and elitist multiobjective genetic algorithm: NSGA-II. IEEE Transactions on Evolutionary Computation \textbf{6}(2), 182--197 (2002)

\bibitem{santos2025}
Santos, G.O., Vieira, L.S., Rossetti, G., Ferreira, C.H.G., Moreira, G.J.P.: A high-performance evolutionary multiobjective community detection algorithm. Social Network Analysis and Mining \textbf{15}(1), 110 (2025). \url{https://doi.org/10.1007/s13278-025-01519-7}

\bibitem{newman2004}
Newman, M.E.J., Girvan, M.: Finding and evaluating community structure in networks. Physical Review E \textbf{69}(2), 026113 (2004)

\bibitem{wen2016}
Wen, X., Chen, W.N., Lin, Y., Gu, T., Pan, S., Li, Y., Yin, J.: A maximal clique based multiobjective evolutionary algorithm for overlapping community detection. IEEE Transactions on Evolutionary Computation \textbf{21}(5), 797--812 (2016)

\bibitem{tian2020}
Tian, D., Zhao, X., Zhao, Z., Shi, Z., Shen, Y.: An evolutionary multiobjective optimization-based fuzzy method for overlapping community detection. IEEE Transactions on Fuzzy Systems \textbf{29}(10), 2843--2856 (2020)

\bibitem{bello2018}
Bello-Orgaz, G., Salcedo-Sanz, S., Camacho, D.: A multiobjective genetic algorithm for overlapping community detection based on edge encoding. Information Sciences \textbf{462}, 290--314 (2018)

\bibitem{xie2011}
Xie, J., Szymanski, B.K., Liu, X.: SLPA: Uncovering Overlapping Communities in Social Networks via A Speaker-listener Interaction Dynamic Process. In: IEEE 11th International Conference on Data Mining Workshops (ICDMW), pp. 344--349 (2011)

\bibitem{ponomarenko2021}
Ponomarenko, A., Dix, A., Muchnik, L.: LPAM: Link partitioning around medoids for overlapping community detection. PLOS ONE \textbf{16}(4), e0248744 (2021)

\bibitem{shchur2019}
Shchur, O., G{\"u}nnemann, S.: Overlapping community detection on complex networks with graph convolutional networks. In: Proceedings of the ACM SIGKDD International Conference on Knowledge Discovery \& Data Mining (2019)

\bibitem{lancichinetti2009onmi}
Lancichinetti, A., Fortunato, S., Kert{\'e}sz, J.: Detecting the overlapping and hierarchical community structure in complex networks. New Journal of Physics \textbf{11}(3), 033015 (2009)

\bibitem{lancichinetti2008lfr}
Lancichinetti, A., Fortunato, S., Radicchi, F.: Benchmark graphs for testing community detection algorithms. Physical Review E \textbf{78}(4), 046110 (2008)

\bibitem{shen2009}
Shen, H., Cheng, X., Cai, K., Hu, M.B.: Detect overlapping and hierarchical community structure in networks. Physica A: Statistical Mechanics and its Applications \textbf{388}(8), 1706--1712 (2009)

\bibitem{nicosia2009}
Nicosia, V., Mangioni, G., Carchiolo, V., Malgeri, M.: Extending the definition of modularity to directed graphs with overlapping communities. Journal of Statistical Mechanics: Theory and Experiment \textbf{2009}(03), P03024 (2009)

\bibitem{zachary1977}
Zachary, W.W.: An information flow model for conflict and fission in small groups. Journal of Anthropological Research \textbf{33}(4), 452--473 (1977). \url{https://doi.org/10.1086/jar.33.4.3629752}

\bibitem{lusseau2003}
Lusseau, D., Schneider, K., Boisseau, O.J., Haase, P., Slooten, E., Dawson, S.M.: The bottlenose dolphin community of Doubtful Sound features a large proportion of long-lasting associations. Behavioral Ecology and Sociobiology \textbf{54}(4), 396--405 (2003). \url{https://doi.org/10.1007/s00265-003-0651-y}

\bibitem{leskovec2007}
Leskovec, J., Kleinberg, J., Faloutsos, C.: Graph evolution: Densification and shrinking diameters. ACM Transactions on Knowledge Discovery from Data \textbf{1}(1), 2--es (2007). \url{https://snap.stanford.edu/data/email-Eu-core.html}

\bibitem{yang2015}
Yang, J., Leskovec, J.: Defining and evaluating network communities based on ground-truth. Knowledge and Information Systems \textbf{42}(1), 181--213 (2015). \url{https://snap.stanford.edu/data/com-DBLP.html}

\end{thebibliography}
}

\end{document}
